\documentclass[a4paper,12pt]{report}

% -----------------------------------------------------------------------
% Rapport de stage / PFA — structure inspirée d'un rapport EMSI, adaptée
% au sujet : « Conception d'une Plateforme de Gestion des Bénéficiaires
% de l'Axe Économie Sociale et Solidaire au sein de la Fondation OCP en
% intégrant des Exigences de Cybersécurité ».
%
% IMPORTANT : ce rapport documente le développement de Secure
% Requirements Studio (l'outil d'ingénierie des exigences) et son
% rapport avec le projet cible (la future plateforme bénéficiaires).
% Il ne prétend pas que la plateforme bénéficiaires a été implémentée.
% Toute information non disponible est marquée [À COMPLÉTER] plutôt
% qu'inventée.
% -----------------------------------------------------------------------

\usepackage[utf8]{inputenc}
\usepackage[T1]{fontenc}
\usepackage[french]{babel}
\usepackage{geometry}
\geometry{a4paper, margin=2.5cm}
\usepackage{graphicx}
\usepackage{hyperref}
\usepackage{listings}
\usepackage{xcolor}
\usepackage{booktabs}
\usepackage{longtable}
\usepackage{fancyhdr}
\usepackage{titlesec}
\usepackage{tocbibind}

\hypersetup{
    colorlinks=true,
    linkcolor=black,
    citecolor=black,
    urlcolor=blue,
    pdftitle={Conception d'une Plateforme de Gestion des Bénéficiaires - Fondation OCP},
}

\lstset{
    basicstyle=\ttfamily\small,
    breaklines=true,
    frame=single,
    backgroundcolor=\color{gray!10},
    extendedchars=true,
    literate=
        {é}{{\'e}}1 {è}{{\`e}}1 {à}{{\`a}}1 {ç}{{\c c}}1
        {â}{{\^a}}1 {î}{{\^i}}1 {ô}{{\^o}}1 {ê}{{\^e}}1
        {û}{{\^u}}1 {ù}{{\`u}}1 {É}{{\'E}}1 {È}{{\`E}}1
        {À}{{\`A}}1 {Ç}{{\c C}}1 {—}{{--}}1,
}

\pagestyle{fancy}
\fancyhf{}
\fancyhead[L]{Rapport de stage / PFA}
\fancyhead[R]{\leftmark}
\fancyfoot[C]{\thepage}

\newcommand{\TODO}[1]{\textcolor{red}{\textbf{[À COMPLÉTER : #1]}}}

\begin{document}

% ------------------------- Page de garde -------------------------------
\begin{titlepage}
    \centering
    \vspace*{1cm}
    {\large \TODO{Nom de l'école / logo EMSI}}\\[1.5cm]
    {\Large \textbf{Rapport de Stage / Projet de Fin d'Année (PFA)}}\\[1cm]
    \rule{\linewidth}{0.5pt}\\[0.5cm]
    {\LARGE \textbf{Conception d'une Plateforme de Gestion des Bénéficiaires\\
    de l'Axe Économie Sociale et Solidaire au sein de la\\
    Fondation OCP en intégrant des Exigences de Cybersécurité}}\\[0.5cm]
    \rule{\linewidth}{0.5pt}\\[1.5cm]

    {\large Réalisé par~: \TODO{Nom(s) de l'étudiant(e)}}\\[0.5cm]
    {\large Encadrant académique~: \TODO{Nom de l'encadrant EMSI}}\\[0.3cm]
    {\large Encadrant professionnel (Fondation OCP)~: \TODO{Nom de l'encadrant OCP}}\\[1.5cm]

    {\large Année universitaire~: \TODO{20XX/20XX}}\\[2cm]

    \vfill
    {\large Fondation OCP — \TODO{Direction / Axe Économie Sociale et Solidaire}}
\end{titlepage}

\pagenumbering{roman}

% ------------------------- Remerciements -------------------------------
\chapter*{Remerciements}
\addcontentsline{toc}{chapter}{Remerciements}
\TODO{Remerciements personnalisés envers l'encadrant académique, l'encadrant professionnel, l'équipe d'accueil, etc.}

% ------------------------- Résumé / Abstract ----------------------------
\chapter*{Résumé}
\addcontentsline{toc}{chapter}{Résumé}
Ce rapport documente la réalisation de \textbf{Secure Requirements Studio (SRS)}, un outil interne
d'ingénierie des exigences conçu pour cadrer, spécifier et sécuriser le projet de
Plateforme de Gestion des Bénéficiaires de l'Axe Économie Sociale et Solidaire de la
Fondation OCP. Le stage s'est concentré sur la conception et le développement de SRS
lui-même — un entretien structuré (assisté par IA ou en questionnaire classique) qui
transforme les réponses des parties prenantes en un Cahier des Charges, une conception
applicative, un diagramme de conception MVP et des exigences de cybersécurité
directement traçables aux réponses collectées. \TODO{Compléter avec les résultats
quantitatifs obtenus et les limites précises constatées.}

\textbf{Mots-clés~:} ingénierie des exigences, cybersécurité, Fondation OCP, économie
sociale et solidaire, cahier des charges, intelligence artificielle.

\chapter*{Abstract}
\addcontentsline{toc}{chapter}{Abstract}
\TODO{Traduction anglaise du résumé.}

\tableofcontents
\listoffigures
\listoftables

\clearpage
\pagenumbering{arabic}

% -------------------------------------------------------------------
\chapter{Introduction générale}

Dans le cadre de \TODO{stage de fin d'année / PFA, nom du diplôme visé, durée exacte du
stage}, ce rapport présente le travail réalisé au sein de la Fondation OCP autour du
projet~: \emph{« Conception d'une Plateforme de Gestion des Bénéficiaires de l'Axe
Économie Sociale et Solidaire au sein de la Fondation OCP en intégrant des Exigences de
Cybersécurité »}.

La Fondation OCP pilote un ensemble de programmes sociaux et économiques regroupés sous
l'Axe Économie Sociale et Solidaire, au bénéfice de coopératives et de bénéficiaires
directs et indirects à travers le Maroc. Avant de développer une plateforme numérique de
gestion pour cet axe, il est indispensable de cadrer précisément les besoins métier, les
contraintes techniques et — étant donné la sensibilité des données traitées — les
exigences de cybersécurité et de protection des données personnelles.

C'est ce besoin de cadrage rigoureux qui a motivé la réalisation de
\textbf{Secure Requirements Studio (SRS)}~: un outil interne d'ingénierie des exigences
qui mène un entretien structuré avec les parties prenantes (en mode assisté par IA ou en
questionnaire classique hors ligne) et transforme progressivement leurs réponses en un
Cahier des Charges, une conception applicative, un diagramme de conception MVP et un
ensemble d'exigences de cybersécurité — chaque élément généré restant traçable à une
réponse effectivement donnée pendant l'entretien, plutôt qu'à un modèle générique
recopié d'un projet à l'autre.

\section{Objectifs du stage}
\TODO{Lister les objectifs pédagogiques et professionnels assignés par l'encadrant.}

\section{Structure du rapport}
Ce rapport est structuré comme suit~: le chapitre~2 présente l'organisme d'accueil ; le
chapitre~3 détaille le contexte et la problématique ; le chapitre~4 couvre l'analyse et
la conception du système ; le chapitre~5 décrit la réalisation de Secure Requirements
Studio ; le chapitre~6 explique le rôle de l'assistance par intelligence artificielle ;
le chapitre~7 détaille les aspects de cybersécurité ; le chapitre~8 présente les tests
effectués ; le chapitre~9 discute les résultats obtenus et les limites ; le chapitre~10
conclut le rapport et ouvre sur les perspectives.

\chapter{Présentation de la Fondation OCP}

\section{Présentation générale}
\TODO{Présentation officielle de la Fondation OCP : date de création, statut juridique,
rattachement au Groupe OCP, gouvernance.}

\section{Mission et domaines d'intervention}
\TODO{Domaines d'intervention de la Fondation (éducation, entrepreneuriat, économie
sociale et solidaire, environnement, culture, etc.), sourcés depuis des documents
officiels ou le site institutionnel de la Fondation.}

\section{Contexte de l'Axe Économie Sociale et Solidaire}
L'Axe Économie Sociale et Solidaire regroupe les programmes de la Fondation destinés aux
coopératives et à leurs bénéficiaires (agriculteurs, femmes, jeunes, personnes en
situation de handicap, actifs et inactifs). Ces programmes s'inscrivent également dans
une logique de suivi par rapport aux Objectifs de Développement Durable (ODD) et aux
indicateurs ESG (Environnement, Social, Gouvernance).

\TODO{Compléter avec le nombre exact de coopératives accompagnées, la répartition
géographique, et les chiffres officiels communiqués par la Direction concernée — ne pas
extrapoler des chiffres non confirmés.}

\chapter{Contexte et problématique}

\section{Contexte du projet}
Le suivi des coopératives et de leurs bénéficiaires au sein de l'Axe Économie Sociale et
Solidaire repose aujourd'hui, selon les informations recueillies, en partie sur des
fichiers bureautiques (Excel, CSV) et des remontées manuelles, ce qui limite la
visibilité consolidée et complique la production de rapports fiables. \TODO{Confirmer
avec la Direction concernée l'état exact de l'existant avant publication du rapport
final — ne pas présenter cette description comme un fait officiel non vérifié.}

\section{Problématique}
Comment concevoir, avant tout développement, une spécification complète, cohérente et
sécurisée de la future plateforme de gestion des bénéficiaires — couvrant les besoins
fonctionnels, les contraintes techniques et les exigences de cybersécurité — à partir
d'un entretien structuré avec les parties prenantes, plutôt que d'un cahier des charges
générique ?

\section{Besoin métier}
La Fondation OCP a besoin d'un cadrage rigoureux, traçable et réutilisable des exigences
du futur système, garantissant que chaque exigence documentée provient effectivement
d'une réponse de partie prenante, d'une règle métier explicite, ou d'une hypothèse
clairement identifiée comme telle — et non d'une invention.

\section{Objectifs du projet}
\begin{itemize}
    \item Fournir un outil d'entretien structuré (IA ou questionnaire classique) dédié à
    ce projet précis, et non un chatbot généraliste.
    \item Générer un Cahier des Charges concis (3 à 5~pages) directement dérivé des
    réponses collectées.
    \item Produire une conception applicative de haut niveau et un diagramme de
    conception MVP unique.
    \item Générer des exigences de cybersécurité proportionnées aux risques
    effectivement identifiés durant l'entretien.
    \item Garantir qu'aucune information confidentielle de la Fondation OCP n'est
    inventée par le système.
\end{itemize}

\section{Périmètre}
Le périmètre de ce stage couvre la conception et le développement de
\textbf{Secure Requirements Studio}, l'outil d'ingénierie des exigences. Il ne couvre
\textbf{pas} le développement de la future Plateforme de Gestion des Bénéficiaires
elle-même, qui reste un projet ultérieur distinct, à charge de l'équipe de
développement qui reprendra les livrables produits par SRS.

\chapter{Analyse et conception}

\section{Étude de l'existant}
\TODO{Décrire les outils/processus existants avant SRS pour le cadrage de projets
similaires au sein de la Fondation OCP, si une étude a été menée. Sinon, indiquer
explicitement qu'aucune étude formelle de l'existant n'a été réalisée dans le cadre de
ce stage.}

\section{Analyse des besoins}
Les besoins ont été structurés autour d'un modèle de spécification centralisé (le
\emph{Knowledge Graph}), qui organise les réponses de l'entretien par domaine~:
parties prenantes, utilisateurs, coopératives, bénéficiaires, exigences fonctionnelles,
règles métier, workflows, documents, KPI, ODD, ESG, sécurité, données, intégrations,
architecture, infrastructure, notifications, reporting, tests et déploiement.

\section{Acteurs et rôles}
\begin{itemize}
    \item \textbf{Administrateur / consultant SRS} — mène l'entretien de cadrage.
    \item \textbf{Partie prenante interrogée} — représentant métier de la Fondation OCP
    ou de l'Axe Économie Sociale et Solidaire.
    \item \textbf{Équipe de développement (destinataire des livrables)} — reprend le
    Cahier des Charges et la conception pour développer la future plateforme.
\end{itemize}

\section{Fonctionnalités principales de SRS}
\begin{itemize}
    \item Entretien structuré (Classic Questionnaire hors ligne ou IA adaptative).
    \item Moteur de Knowledge Graph avec suivi de complétude et de dépendances entre
    concepts.
    \item Moteur de génération d'exigences (fonctionnelles, non fonctionnelles,
    sécurité, techniques).
    \item Moteur d'analyse de sécurité (score de sécurité, checklist, matrice RBAC,
    registre des risques, analyse d'impact vie privée).
    \item Moteur de revue (règles de complétude avant export final).
    \item Générateur de documents multi-formats (JSON, Markdown, HTML, DOCX, PDF).
\end{itemize}

\section{Architecture générale}
SRS suit une architecture classique en trois couches~:

\begin{itemize}
    \item \textbf{Frontend} — application Next.js (React) assurant l'assistant
    d'entretien, les tableaux de bord et la consultation des livrables générés.
    \item \textbf{Backend} — API FastAPI (Python) exposant les routes projet, entretien,
    knowledge graph, exigences, sécurité, revue et documents.
    \item \textbf{Base de données} — PostgreSQL en production (SQLite en environnement
    de test), via SQLAlchemy asynchrone.
\end{itemize}

Il n'existe volontairement \textbf{aucune couche d'authentification} dans SRS~: chaque
projet appartient à un contexte interne unique (« Fondation OCP »), l'outil étant destiné
à un usage interne restreint pendant la phase de cadrage plutôt qu'à un accès multi-
organisation.

\section{Conception de la solution}
La conception distingue explicitement, à chaque étape de la génération documentaire~:
les informations \emph{confirmées} par une réponse d'entretien, les
\emph{hypothèses} formulées par le système, et les points \emph{à confirmer} avec les
parties prenantes concernées (par exemple les exigences légales/RGPD ou Loi~09-08, qui
sont systématiquement étiquetées comme à valider avec le responsable juridique/DPO).

\section{Diagramme de conception du MVP}
Conformément au principe « un seul diagramme de conception MVP compréhensible par des
publics techniques et non techniques », SRS génère dynamiquement un diagramme unique
(format Mermaid) représentant les acteurs, le frontend, le backend, les modules
effectivement justifiés par les exigences confirmées, la base de données, et les
contrôles de sécurité applicables — sans ajout systématique de briques non justifiées
(microservices, Kubernetes, files de messages, etc.).

\begin{figure}[h]
    \centering
    \TODO{Insérer une capture d'écran ou un rendu du diagramme MVP généré pour un
    projet réel, une fois disponible.}
    \caption{Exemple de diagramme de conception MVP généré par SRS}
\end{figure}

\section{Flux principaux}
\subsection{Flux d'entretien}
Démarrage de session → question adaptative ou issue du catalogue de concepts →
interprétation de la réponse → mise à jour du Knowledge Graph → génération incrémentale
des exigences associées → question suivante ou fin d'entretien.

\subsection{Flux de génération documentaire}
Sur demande, le générateur de documents assemble le contenu à partir de l'état courant
du Knowledge Graph, des exigences et — si disponible — de l'analyse de sécurité, puis
l'exporte dans le format demandé (JSON, Markdown, HTML, DOCX, PDF).

\subsection{Flux de revue}
Avant export final, le moteur de revue vérifie la complétude (parties prenantes,
exigences, sécurité, architecture, tests) et distingue les éléments confirmés, supposés
et restant à confirmer.

\section{Choix technologiques}
\begin{longtable}{|l|l|}
\hline
\textbf{Composant} & \textbf{Technologie} \\
\hline
Backend & FastAPI, SQLAlchemy (async), Alembic \\
Base de données & PostgreSQL (prod), SQLite (tests) \\
Frontend & Next.js, React, Tailwind CSS \\
IA (entretien adaptatif) & Google Gemini, via une abstraction de client IA \\
Export documentaire & python-docx, ReportLab, markdown2 \\
Tests backend & pytest, pytest-asyncio \\
\hline
\end{longtable}

\chapter{Réalisation de Secure Requirements Studio}

\section{Environnement de développement}
\begin{itemize}
    \item Backend~: Python 3.12, FastAPI, Poetry pour la gestion des dépendances.
    \item Frontend~: Node.js, Next.js 15, TypeScript.
    \item Éditeur~: \TODO{VS Code / autre, selon l'environnement réellement utilisé}.
    \item Gestion de version~: Git. \TODO{Préciser l'hébergement du dépôt (interne,
    GitHub, GitLab...) et la politique de branches si applicable.}
\end{itemize}

\section{Architecture de Secure Requirements Studio}
Le backend est organisé en couches~: \texttt{routers} (endpoints HTTP), \texttt{services}
(logique métier des différents moteurs), \texttt{repositories} (accès aux données),
\texttt{models} (entités SQLAlchemy) et \texttt{domain} (connaissances métier statiques,
telles que le catalogue de concepts et les référentiels de sécurité).

\section{Questionnaire guidé}
Le \emph{Classic Questionnaire} permet de mener un entretien complet sans dépendre d'un
fournisseur d'IA~: chaque question est associée à un domaine du Knowledge Graph, avec
suivi automatique de la complétude et des dépendances entre concepts (par exemple, une
question sur le RBAC régional n'est posée qu'après confirmation du regroupement des
coopératives par région).

\section{Entretien assisté par IA}
En mode assisté, l'IA joue le rôle d'un consultant senior en ingénierie des exigences~:
questions courtes et précises en français, pas de répétition d'une question déjà posée
sous la même forme, et restriction stricte au périmètre du projet cadré. Toute question
hors périmètre reçoit une réponse de recentrage plutôt qu'une réponse de chatbot
généraliste.

\section{Intégration des fournisseurs IA}
\TODO{Décrire précisément, au moment de la rédaction finale, les fournisseurs IA
effectivement intégrés et testés (Gemini confirmé dans le code ; les autres
fournisseurs mentionnés dans les spécifications initiales — OpenAI, Anthropic Claude,
Kimi, Ollama — ne doivent être décrits comme livrés que s'ils ont réellement été
implémentés et testés au moment de la rédaction).}

\section{Génération dynamique des exigences}
Chaque réponse suffisamment complète déclenche la génération d'exigences typées
(fonctionnelle, non fonctionnelle, sécurité, technique, etc.), chacune associée à une
clé unique, un identifiant du concept source (\texttt{source\_concept\_key}), une
priorité et un statut — garantissant la traçabilité entre une exigence et la réponse
qui l'a produite.

\section{Génération du Cahier des Charges}
Le Cahier des Charges est assemblé automatiquement à partir de l'état courant du
Knowledge Graph et des exigences générées, avec pour contrainte de rester concis
(3 à 5~pages) plutôt que de reproduire un document générique volumineux.

\section{Génération de la conception}
La conception applicative de haut niveau (architecture, acteurs, modules,
flux principaux) est dérivée du même modèle de données que le Cahier des Charges,
garantissant la cohérence entre les deux documents.

\section{Génération du diagramme MVP}
Voir section~4.7 pour le détail du principe de génération du diagramme unique de
conception MVP.

\section{Génération des documents}
Les formats d'export pris en charge sont~: JSON (données machine), Markdown (source de
vérité textuelle), HTML (dérivé du Markdown), DOCX et PDF (générés directement à partir
du modèle de contenu structuré, pour garantir une mise en forme fiable indépendamment de
la conversion Markdown).

\chapter{Intelligence artificielle et assistance à l'analyse}

\section{Objectif de l'assistant IA}
L'assistant IA n'est pas le produit final de SRS~: c'est un intervieweur intelligent
dont le rôle est de collecter des informations. Le produit réel est la transformation
des réponses de l'entretien en modèle de projet structuré, puis en exigences, en
architecture et en documents.

\section{Fonctionnement de l'entretien intelligent}
À chaque tour, l'IA reçoit un contexte structuré (concept courant, historique des
échanges, informations déjà connues, informations manquantes) plutôt qu'une simple
conversation non structurée, afin de générer une question courte, précise et non
redondante.

\section{Gestion du contexte}
Le contexte transmis à l'IA distingue explicitement les informations connues du projet,
les hypothèses, et les questions nécessitant une confirmation — évitant que l'IA
présente une supposition comme un fait acquis.

\section{Restrictions du domaine de l'IA}
L'IA est strictement cantonnée aux sujets liés au cadrage de la plateforme (besoins,
architecture, données, sécurité, déploiement, tests). Toute question hors périmètre
reçoit un message de recentrage explicite plutôt qu'une réponse de chatbot généraliste.

\section{Gestion des fournisseurs IA}
Le client IA est conçu derrière une interface (\texttt{AIClient}) découplée du reste de
l'application, afin de pouvoir ajouter ou remplacer un fournisseur sans modifier le
moteur d'entretien lui-même.

\section{Fournisseurs IA intégrés}
\TODO{Lister ici, avec exactitude, l'état réel de chaque fournisseur au moment de la
rédaction finale (implémenté et testé / partiellement implémenté / non implémenté),
plutôt que l'ambition initiale du cahier des charges.}

\section{Sécurité des clés API}
Les clés API des fournisseurs IA ne transitent jamais côté client~: elles sont lues
côté serveur à partir de variables d'environnement, ne sont jamais journalisées, ni
incluses dans les documents générés, ni committées dans le dépôt Git.

\section{Limites et risques liés à l'IA}
\begin{itemize}
    \item Une réponse d'IA reste une interprétation~: elle doit toujours être
    distinguée d'une information officiellement confirmée par une partie prenante.
    \item En cas d'indisponibilité du fournisseur IA, le Classic Questionnaire reste
    pleinement fonctionnel afin de ne jamais bloquer un entretien.
    \item Aucune information confidentielle non confirmée par la Fondation OCP ne doit
    être présentée comme un fait acquis dans les documents générés.
\end{itemize}

\chapter{Cybersécurité}

La cybersécurité est un axe central du projet, tant dans l'outil SRS lui-même que dans
les exigences qu'il génère pour la future plateforme.

\section{Modèle de menace}
Le moteur de sécurité de SRS applique une analyse de type \textbf{STRIDE} (Spoofing,
Tampering, Repudiation, Information disclosure, Denial of service, Elevation of
privilege), pondérée selon une logique \textbf{DREAD}, sur un catalogue de concepts
prioritaires du Knowledge Graph (authentification, autorisation, protection des
données, disponibilité, etc.).

\section{Analyse des risques}
Pour chaque risque identifié, un registre structuré est généré, incluant~: identifiant,
actif concerné, menace, vulnérabilité, impact, probabilité, niveau de risque,
mitigation proposée et risque résiduel. Seuls les risques dont le niveau est
\emph{élevé} ou \emph{critique} apparaissent dans le registre final, afin de ne pas
diluer l'attention sur des risques mineurs.

\section{Contrôle d'accès}
Une matrice RBAC (Role-Based Access Control) est générée à partir des acteurs
effectivement identifiés dans les exigences fonctionnelles — elle n'est donc pas fixe,
mais reflète les rôles confirmés durant l'entretien (par exemple~: administrateur,
coordinateur régional, représentant de coopérative).

\section{Protection des données}
Une analyse d'impact relative à la protection des données (PIA) est générée, avec une
règle stricte~: les exigences légales telles que la Loi~09-08 ou le RGPD ne sont
\textbf{jamais} présentées comme confirmées par défaut. Elles sont systématiquement
étiquetées comme « à confirmer avec le responsable juridique / DPO », sauf confirmation
explicite d'une partie prenante durant l'entretien.

\section{Sécurité des fichiers}
\TODO{Détailler les contrôles de sécurité relatifs au téléversement de fichiers
(validation de type, taille, scan antivirus) une fois ces exigences confirmées pour un
projet concret — ne pas décrire de contrôles qui ne sont pas encore documentés.}

\section{Journalisation et traçabilité}
Chaque exigence générée référence la réponse d'entretien dont elle est issue
(\texttt{source\_concept\_key}), fournissant une traçabilité de bout en bout entre une
exigence du Cahier des Charges et la réponse de la partie prenante qui l'a motivée.

\section{Sauvegarde et continuité}
\TODO{Documenter la stratégie de sauvegarde et de continuité effectivement retenue pour
l'environnement d'hébergement de SRS, une fois le déploiement finalisé.}

\section{Secure SDLC}
Le développement de SRS applique des pratiques de cycle de développement sécurisé de
base~: séparation stricte de la configuration sensible (clés API) via variables
d'environnement, absence de journalisation des secrets, et couverture par tests
automatisés des principaux flux applicatifs (voir chapitre~8).

\chapter{Tests}

\section{Tests fonctionnels}
Le backend est couvert par une suite de tests automatisés (\texttt{pytest},
\texttt{pytest-asyncio}) exécutée contre une base SQLite en mémoire, couvrant notamment~:
la création et le cycle de vie des projets, l'amorçage automatique du Knowledge Graph,
le déroulement complet d'un entretien, la génération des exigences, l'analyse de
sécurité, le moteur de revue et la génération documentaire multi-formats.

\TODO{Indiquer le nombre exact de tests et le taux de couverture au moment de la
rédaction finale du rapport (à régénérer via \texttt{pytest --cov} juste avant la
remise).}

\section{Tests d'intégration}
Les tests d'API exercent les endpoints HTTP de bout en bout via un client de test
asynchrone, garantissant que les routes, les schémas Pydantic et les services métier
restent cohérents entre eux.

\section{Tests de sécurité}
Le moteur de sécurité est testé sur des cas explicites~: absence de fabrication d'une
base légale de traitement des données non confirmée, cohérence entre les acteurs
déclarés dans les exigences et la matrice RBAC générée, et non-duplication d'une
analyse de sécurité déjà exécutée pour un même projet.

\section{Tests de génération documentaire}
Le générateur de documents est testé pour garantir qu'un projet fraîchement créé (sans
réponse d'entretien) ne produit \textbf{aucune donnée fabriquée}~: dictionnaire de
données vide, absence d'analyse de sécurité tant qu'elle n'a pas été exécutée, et
mention explicite des sections volontairement non générées (diagrammes BPMN, maquettes
visuelles).

\section{Tests de l'assistant IA}
\TODO{Décrire les tests spécifiques à l'assistant IA (longueur des questions, respect
du périmètre du projet, gestion des erreurs de connexion au fournisseur) une fois
formalisés et automatisés.}

\chapter{Résultats}

\section{Fonctionnalités réalisées}
Au moment de la rédaction, Secure Requirements Studio permet~:
\begin{itemize}
    \item de mener un entretien de cadrage complet, en mode Classic Questionnaire ou
    assisté par IA (Gemini) ;
    \item de générer dynamiquement des exigences fonctionnelles, non fonctionnelles, de
    sécurité et techniques, traçables aux réponses de l'entretien ;
    \item de produire une analyse de sécurité (score, checklist, matrice RBAC, registre
    des risques, analyse d'impact vie privée) ;
    \item d'exécuter une revue de complétude avant export ;
    \item de générer un Cahier des Charges et une conception applicative avec un
    diagramme MVP unique, exportables en JSON, Markdown, HTML, DOCX et PDF.
\end{itemize}

\section{Résultats obtenus}
\TODO{Compléter avec des résultats concrets une fois un ou plusieurs entretiens réels
menés avec des parties prenantes de la Fondation OCP — ne pas anticiper de résultats
métier qui n'ont pas encore été produits.}

\section{Limites}
\begin{itemize}
    \item SRS produit une \emph{conception et une spécification}, pas une implémentation
    de la future Plateforme de Gestion des Bénéficiaires elle-même.
    \item L'intégration de certains fournisseurs IA envisagés initialement (OpenAI,
    Anthropic Claude, Kimi, Ollama) peut être partielle ou absente selon l'état réel du
    code au moment de la rédaction — voir chapitre~6.
    \item L'outil ne dispose volontairement d'aucune authentification, ce qui limite
    son usage à un contexte interne restreint.
    \item \TODO{Ajouter toute autre limite constatée en conditions réelles d'usage.}
\end{itemize}

\section{Perspectives}
\begin{itemize}
    \item Étendre la couverture des fournisseurs IA supportés.
    \item Enrichir le catalogue de concepts du Knowledge Graph au fur et à mesure des
    entretiens réels menés avec les parties prenantes de l'Axe Économie Sociale et
    Solidaire.
    \item Transmettre les livrables générés (Cahier des Charges, conception, diagramme
    MVP, exigences de sécurité) à l'équipe de développement chargée de la future
    Plateforme de Gestion des Bénéficiaires.
\end{itemize}

\chapter{Conclusion générale}

Ce stage a porté sur la conception et le développement de \textbf{Secure Requirements
Studio}, un outil interne d'ingénierie des exigences destiné à cadrer le futur projet de
Plateforme de Gestion des Bénéficiaires de l'Axe Économie Sociale et Solidaire de la
Fondation OCP, en intégrant dès l'origine des exigences de cybersécurité.

Le travail réalisé a permis de mettre en place un entretien structuré — assisté par IA
ou en questionnaire classique hors ligne — capable de transformer les réponses des
parties prenantes en un modèle de projet structuré (Knowledge Graph), puis en exigences
traçables, en une analyse de sécurité proportionnée aux risques réels, et en un Cahier
des Charges concis accompagné d'une conception applicative et d'un diagramme de
conception MVP unique.

Au-delà de l'aspect technique, ce stage a été l'occasion de travailler un principe
directeur exigeant~: ne jamais présenter une supposition ou une inférence de l'IA comme
un fait acquis de la Fondation OCP, et distinguer systématiquement ce qui est confirmé,
supposé, ou à confirmer. Ce principe a structuré aussi bien la conception du moteur de
génération d'exigences que celle du moteur de sécurité.

\TODO{Compléter cette conclusion avec un bilan personnel du stage (apprentissages
techniques et humains, difficultés rencontrées et comment elles ont été surmontées) une
fois le stage arrivé à son terme.}

Les livrables produits par Secure Requirements Studio — Cahier des Charges, conception
applicative, diagramme MVP et exigences de cybersécurité — constituent la base sur
laquelle l'équipe de développement pourra s'appuyer pour concevoir, dans un second
temps, la Plateforme de Gestion des Bénéficiaires elle-même.


% -------------------------------------------------------------------
\appendix
\chapter{Glossaire}

\begin{description}
    \item[Cahier des Charges] Document synthétique (3 à 5~pages dans SRS) décrivant les
    besoins fonctionnels et non fonctionnels essentiels d'un projet.
    \item[Knowledge Graph] Modèle interne de SRS structurant les réponses de l'entretien
    par domaine et par dépendances entre concepts.
    \item[RBAC] \emph{Role-Based Access Control} — contrôle d'accès basé sur les rôles.
    \item[STRIDE] Méthode de modélisation des menaces (Spoofing, Tampering,
    Repudiation, Information disclosure, Denial of service, Elevation of privilege).
    \item[DREAD] Méthode de pondération des risques (Damage, Reproducibility,
    Exploitability, Affected users, Discoverability).
    \item[PIA] \emph{Privacy Impact Assessment} — analyse d'impact relative à la
    protection des données.
    \item[MVP] \emph{Minimum Viable Product} — version minimale mais fonctionnelle d'un
    produit.
    \item[ODD] Objectifs de Développement Durable (Nations Unies).
    \item[ESG] Environnement, Social, Gouvernance — cadre d'indicateurs de durabilité.
    \item[SRS] Secure Requirements Studio — l'outil objet de ce rapport.
\end{description}

\chapter{Bibliographie et Webographie}

Les références utilisées sont listées dans le fichier \texttt{references.bib} et citées
au fil du texte selon le style bibliographique standard. \TODO{Ajouter les références
effectivement consultées durant le stage (documentation FastAPI, Next.js, OWASP ASVS,
OWASP Top 10, documentation Fondation OCP, etc.) au fur et à mesure de la rédaction.}

\begin{itemize}
    \item OWASP Foundation, \emph{OWASP Top 10}, \url{https://owasp.org/www-project-top-ten/}.
    \item OWASP Foundation, \emph{OWASP Application Security Verification Standard
    (ASVS)}, \url{https://owasp.org/www-project-application-security-verification-standard/}.
    \item FastAPI, documentation officielle, \url{https://fastapi.tiangolo.com/}.
    \item Next.js, documentation officielle, \url{https://nextjs.org/docs}.
    \item \TODO{Site institutionnel de la Fondation OCP (page Économie Sociale et
    Solidaire) — ajouter l'URL exacte consultée.}
\end{itemize}

\chapter{Guide d'installation}

\section{Prérequis}
\begin{itemize}
    \item Python~3.12 et Poetry (backend).
    \item Node.js (frontend Next.js).
    \item Docker et Docker Compose (base de données PostgreSQL et Redis).
\end{itemize}

\section{Démarrage rapide}
Le dépôt fournit des scripts de démarrage prêts à l'emploi~:
\begin{itemize}
    \item \texttt{run.bat} / \texttt{run.ps1} — Windows.
    \item \texttt{run} / \texttt{runlinux} — Linux / macOS / WSL.
\end{itemize}

Ces scripts vérifient la présence d'un fichier \texttt{backend/.env} (le créent à partir
de \texttt{.env.example} si absent), démarrent les services d'infrastructure via
\texttt{docker-compose.yml} (PostgreSQL, Redis), lancent le backend FastAPI puis le
frontend Next.js, et affichent les URLs locales.

\section{Configuration}
Les variables d'environnement sensibles (clés API des fournisseurs IA, chaîne de
connexion à la base de données) sont définies dans \texttt{backend/.env}, jamais
committées dans le dépôt Git — seul \texttt{backend/.env.example} y figure, sans valeur
réelle.

\section{Lancement manuel (sans script)}
\begin{lstlisting}[language=bash]
# Infrastructure (Postgres + Redis)
docker compose up -d

# Backend
cd backend
poetry install
poetry run uvicorn app.main:app --reload

# Frontend (dans un autre terminal)
cd frontend
npm install
npm run dev
\end{lstlisting}

\section{Tests}
\begin{lstlisting}[language=bash]
# Backend
cd backend
poetry run pytest

# Frontend
cd frontend
npx tsc --noEmit
npm run build
\end{lstlisting}

\TODO{Mettre à jour cette annexe si la procédure de déploiement change avant la remise
finale du rapport (voir également DEPLOYMENT.md à la racine du dépôt, si présent).}

\chapter{Annexes techniques}

\section{Arborescence simplifiée du backend}
\begin{lstlisting}
backend/app/
  routers/       endpoints HTTP (projects, interview, knowledge-graph,
                 requirements, security, review, documents, ai-settings)
  services/      logique métier (interview, knowledge_graph, requirement,
                 security_engine, review_engine, document_generator,
                 document_rendering, llm_client, ai_settings)
  repositories/  accès aux données (une classe par entité)
  models/        entités SQLAlchemy
  schemas/       schémas Pydantic (validation / sérialisation API)
  domain/        connaissances métier statiques (catalogue de concepts,
                 types d'exigences, référentiels de sécurité, savoir FOCP)
\end{lstlisting}

\section{Exemple d'exigence générée}
Extrait illustratif de la structure d'une exigence générée par SRS (les valeurs sont
indicatives)~:

\begin{lstlisting}
{
  "requirement_key": "SEC-VAL-001",
  "requirement_type": "security",
  "title": "Validation des donnees par un role autorise",
  "priority": "haute",
  "source_concept_key": "workflows.monthly_validation",
  "acceptance_criteria": [
    "Seuls les roles autorises peuvent valider une soumission mensuelle."
  ]
}
\end{lstlisting}

\section{Exemple de diagramme MVP (Mermaid)}
Extrait illustratif du diagramme de conception MVP généré dynamiquement (structure
simplifiée — le contenu réel dépend des exigences confirmées pour chaque projet)~:

\begin{lstlisting}
flowchart TD
    A0["Administrateur"] --> FE["Frontend"]
    A1["Coordinateur regional"] --> FE
    FE --> API["Backend / API"]
    API --> M0["Exigences fonctionnelles"]
    M0 --> DB[(Base de donnees)]
    API --> SEC["Controles de securite (RBAC, audit, chiffrement)"]
\end{lstlisting}

\TODO{Remplacer ces extraits par des exemples réels générés pour le projet Fondation
OCP, une fois un entretien complet mené et un cahier des charges effectivement exporté.}


\end{document}
